
\IfFileExists{stacks-project.cls}{%
\documentclass{stacks-project}
}{%
\documentclass{amsart}
}
\usepackage{amssymb}

% For dealing with references we use the comment environment
\usepackage{verbatim}
\newenvironment{reference}{\comment}{\endcomment}
%\newenvironment{reference}{}{}
\newenvironment{slogan}{\comment}{\endcomment}
\newenvironment{history}{\comment}{\endcomment}

% For commutative diagrams we use Xy-pic
\usepackage[all]{xy}

% We use 2cell for 2-commutative diagrams.
\xyoption{2cell}
\UseAllTwocells

% We use multicol for the list of chapters between chapters
\usepackage{multicol}

% This is generally recommended for better output
\usepackage{lmodern}
\usepackage[T1]{fontenc}
\usepackage{textalpha}

% For cross-file-references
\usepackage{xr-hyper}
\newcommand{\maybeexternaldocument}[2]{%
  \IfFileExists{#2.aux}{%
    \externaldocument[#1]{#2}%
  }{%
    \PackageWarning{stack}{Missing #2.aux; external references with prefix #1 unavailable}%
  }%
}
\let\externaldocumentorig\externaldocument
\renewcommand{\externaldocument}[2][]{%
  \IfFileExists{#2.aux}{%
    \externaldocumentorig[#1]{#2}%
  }{%
    \PackageWarning{stack}{Missing #2.aux; external references with prefix #1 unavailable}%
  }%
}

% Package for hypertext links:
\usepackage{hyperref}

% For any local file, say "hello.tex" you want to link to please
% use \externaldocument[hello-]{hello}
\externaldocument[introduction-]{introduction}
\externaldocument[conventions-]{conventions}
\externaldocument[sets-]{sets}
\externaldocument[categories-]{categories}
\externaldocument[topology-]{topology}
\externaldocument[sheaves-]{sheaves}
\externaldocument[sites-]{sites}
\externaldocument[stacks-]{stacks}
\externaldocument[fields-]{fields}
\externaldocument[algebra-]{algebra}
\externaldocument[brauer-]{brauer}
\externaldocument[homology-]{homology}
\externaldocument[derived-]{derived}
\externaldocument[simplicial-]{simplicial}
\externaldocument[more-algebra-]{more-algebra}
\externaldocument[smoothing-]{smoothing}
\externaldocument[modules-]{modules}
\externaldocument[sites-modules-]{sites-modules}
\externaldocument[injectives-]{injectives}
\externaldocument[cohomology-]{cohomology}
\externaldocument[sites-cohomology-]{sites-cohomology}
\externaldocument[dga-]{dga}
\externaldocument[dpa-]{dpa}
\externaldocument[sdga-]{sdga}
\externaldocument[hypercovering-]{hypercovering}
\externaldocument[schemes-]{schemes}
\externaldocument[constructions-]{constructions}
\externaldocument[properties-]{properties}
\externaldocument[morphisms-]{morphisms}
\externaldocument[coherent-]{coherent}
\externaldocument[divisors-]{divisors}
\externaldocument[limits-]{limits}
\externaldocument[varieties-]{varieties}
\externaldocument[topologies-]{topologies}
\externaldocument[descent-]{descent}
\externaldocument[perfect-]{perfect}
\externaldocument[more-morphisms-]{more-morphisms}
\externaldocument[flat-]{flat}
\externaldocument[groupoids-]{groupoids}
\externaldocument[more-groupoids-]{more-groupoids}
\externaldocument[etale-]{etale}
\externaldocument[chow-]{chow}
\externaldocument[intersection-]{intersection}
\externaldocument[pic-]{pic}
\externaldocument[weil-]{weil}
\externaldocument[adequate-]{adequate}
\externaldocument[dualizing-]{dualizing}
\externaldocument[duality-]{duality}
\externaldocument[discriminant-]{discriminant}
\externaldocument[derham-]{derham}
\externaldocument[local-cohomology-]{local-cohomology}
\externaldocument[algebraization-]{algebraization}
\externaldocument[curves-]{curves}
\externaldocument[resolve-]{resolve}
\externaldocument[models-]{models}
\externaldocument[functors-]{functors}
\externaldocument[equiv-]{equiv}
\externaldocument[pione-]{pione}
\externaldocument[etale-cohomology-]{etale-cohomology}
\externaldocument[proetale-]{proetale}
\externaldocument[relative-cycles-]{relative-cycles}
\externaldocument[more-etale-]{more-etale}
\externaldocument[trace-]{trace}
\externaldocument[crystalline-]{crystalline}
\externaldocument[spaces-]{spaces}
\externaldocument[spaces-properties-]{spaces-properties}
\externaldocument[spaces-morphisms-]{spaces-morphisms}
\externaldocument[decent-spaces-]{decent-spaces}
\externaldocument[spaces-cohomology-]{spaces-cohomology}
\externaldocument[spaces-limits-]{spaces-limits}
\externaldocument[spaces-divisors-]{spaces-divisors}
\externaldocument[spaces-over-fields-]{spaces-over-fields}
\externaldocument[spaces-topologies-]{spaces-topologies}
\externaldocument[spaces-descent-]{spaces-descent}
\externaldocument[spaces-perfect-]{spaces-perfect}
\externaldocument[spaces-more-morphisms-]{spaces-more-morphisms}
\externaldocument[spaces-flat-]{spaces-flat}
\externaldocument[spaces-groupoids-]{spaces-groupoids}
\externaldocument[spaces-more-groupoids-]{spaces-more-groupoids}
\externaldocument[bootstrap-]{bootstrap}
\externaldocument[spaces-pushouts-]{spaces-pushouts}
\externaldocument[spaces-chow-]{spaces-chow}
\externaldocument[groupoids-quotients-]{groupoids-quotients}
\externaldocument[spaces-more-cohomology-]{spaces-more-cohomology}
\externaldocument[spaces-simplicial-]{spaces-simplicial}
\externaldocument[spaces-duality-]{spaces-duality}
\externaldocument[formal-spaces-]{formal-spaces}
\externaldocument[restricted-]{restricted}
\externaldocument[spaces-resolve-]{spaces-resolve}
\externaldocument[formal-defos-]{formal-defos}
\externaldocument[defos-]{defos}
\externaldocument[cotangent-]{cotangent}
\externaldocument[examples-defos-]{examples-defos}
\externaldocument[algebraic-]{algebraic}
\externaldocument[examples-stacks-]{examples-stacks}
\externaldocument[stacks-sheaves-]{stacks-sheaves}
\externaldocument[criteria-]{criteria}
\externaldocument[artin-]{artin}
\externaldocument[quot-]{quot}
\externaldocument[stacks-properties-]{stacks-properties}
\externaldocument[stacks-morphisms-]{stacks-morphisms}
\externaldocument[stacks-limits-]{stacks-limits}
\externaldocument[stacks-cohomology-]{stacks-cohomology}
\externaldocument[stacks-perfect-]{stacks-perfect}
\externaldocument[stacks-introduction-]{stacks-introduction}
\externaldocument[stacks-more-morphisms-]{stacks-more-morphisms}
\externaldocument[stacks-geometry-]{stacks-geometry}
\externaldocument[moduli-]{moduli}
\externaldocument[moduli-curves-]{moduli-curves}
\externaldocument[examples-]{examples}
\externaldocument[exercises-]{exercises}
\externaldocument[guide-]{guide}
\externaldocument[desirables-]{desirables}
\externaldocument[coding-]{coding}
\externaldocument[obsolete-]{obsolete}
\externaldocument[fdl-]{fdl}
\externaldocument[index-]{index}

% Heap Project documents
\externaldocument[linear-algebra-]{linear-algebra}
\externaldocument[groups-]{groups}
\externaldocument[complex-analysis-]{complex-analysis}
\externaldocument[functional-analysis-]{functional-analysis}
\externaldocument[categories-]{categories}
\externaldocument[commutative-algebra-]{commutative-algebra}
\externaldocument[ringed-spaces-]{ringed-spaces}
\externaldocument[homological-algebra-]{homological-algebra}
\externaldocument[lie-algebras-]{lie-algebras}
\externaldocument[representation-theory-]{representation-theory}
\externaldocument[smooth-manifolds-]{smooth-manifolds}
\externaldocument[differential-topology-]{differential-topology}
\externaldocument[algebraic-topology-]{algebraic-topology}
\externaldocument[differential-geometry-]{differential-geometry}
\externaldocument[algebraic-geometry-]{algebraic-geometry}
\externaldocument[symplectic-geometry-]{symplectic-geometry}
\externaldocument[complex-geometry-]{complex-geometry}
\externaldocument[hodge-theory-]{hodge-theory}
\maybeexternaldocument{deformations-}{deformations}
\externaldocument[vertex-algebras-]{vertex-algebras}
\externaldocument[springer-]{springer}
\externaldocument[infty-categories-]{infty-categories}
\externaldocument[seminars-differential-geometry-]{%
seminars-differential-geometry}
\externaldocument[seminars-complex-geometry-]{%
seminars-complex-geometry}
\externaldocument[seminars-algebraic-geometry-]{%
seminars-algebraic-geometry}
\externaldocument[seminars-others-]{%
seminars-others}
%\externaldocument[geometric-prequantization-]{geometric-prequantization}
\externaldocument[surfaces-]{surfaces}
\externaldocument[birational-maps-commutative-algebra-]{birational-maps-commutative-algebra}
\externaldocument[cremona-transformations-]{cremona-transformations}
\externaldocument[derived-categories-of-sheaves-]{derived-categories-of-sheaves}
\externaldocument[geometric-invariant-theory-]{geometric-invariant-theory}
\externaldocument[geometric-stability-conditions-and-group-actions-]{geometric-stability-conditions-and-group-actions}
\externaldocument[moduli-spaces-of-sheaves-]{moduli-spaces-of-sheaves}
\externaldocument[bridgeland-stability-general-theory-]{bridgeland-stability-general-theory}
\externaldocument[hilbert-schemes-]{hilbert-schemes}
\externaldocument[noncommutative-abelian-surfaces-and-kummer-type-hyperkahler-varieties-]{noncommutative-abelian-surfaces-and-kummer-type-hyperkahler-varieties}
\externaldocument[mumford-tate-groups-in-hodge-theory-]{mumford-tate-groups-in-hodge-theory}
\externaldocument[hodge-birational-atoms-2026-]{hodge-birational-atoms-2026}
\externaldocument[xxii-egd-2026-]{%
xxii-egd-2026}
\externaldocument[pde-]{pde}
\externaldocument[probability-]{probability}

\externaldocument[linguistics-]{linguistics}
\externaldocument[genealogia-familiar-]{%
genealogia-familiar}

\externaldocument[%
16th-alga-meeting-2026-]{%
16th-alga-meeting-2026}

\externaldocument[%
iv-jornadas-lefschetz-2026-]{%
iv-jornadas-lefschetz-2026}

\externaldocument[reading-the-masters-]{%
reading-the-masters}

\externaldocument[real-analysis-]{%
real-analysis}

\externaldocument[optimal-transport-]{%
optimal-transport}

\externaldocument[history-]{history}

% Theorem environments.
%
\theoremstyle{plain}
\newtheorem{theorem}[subsection]{Theorem}
\newtheorem{proposition}[subsection]{Proposition}
\newtheorem{lemma}[subsection]{Lemma}

\theoremstyle{definition}
\newtheorem{definition}[subsection]{Definition}
\newtheorem{example}[subsection]{Example}
\newtheorem{exercise}[subsection]{Exercise}
\newtheorem{situation}[subsection]{Situation}

\theoremstyle{remark}
\newtheorem{remark}[subsection]{Remark}
\newtheorem{remarks}[subsection]{Remarks}

\numberwithin{equation}{subsection}

% Macros
%
\def\lim{\mathop{\mathrm{lim}}\nolimits}
\def\colim{\mathop{\mathrm{colim}}\nolimits}
\def\Spec{\mathop{\mathrm{Spec}}}
\def\Hom{\mathop{\mathrm{Hom}}\nolimits}
\def\Ext{\mathop{\mathrm{Ext}}\nolimits}
\def\SheafHom{\mathop{\mathcal{H}\!\mathit{om}}\nolimits}
\def\SheafExt{\mathop{\mathcal{E}\!\mathit{xt}}\nolimits}
\def\Sch{\mathit{Sch}}
\def\Mor{\mathop{\mathrm{Mor}}\nolimits}
\def\Ob{\mathop{\mathrm{Ob}}\nolimits}
\def\Sh{\mathop{\mathit{Sh}}\nolimits}
\def\NL{\mathop{N\!L}\nolimits}
\def\CH{\mathop{\mathrm{CH}}\nolimits}
\def\proetale{{pro\text{-}\acute{e}tale}}
\def\etale{{\acute{e}tale}}
\def\QCoh{\mathit{QCoh}}
\def\Ker{\mathop{\mathrm{Ker}}}
\def\Im{\mathop{\mathrm{Im}}}
\def\Coker{\mathop{\mathrm{Coker}}}
\def\Coim{\mathop{\mathrm{Coim}}}

% Boxtimes
%
\DeclareMathSymbol{\boxtimes}{\mathbin}{AMSa}{"02}

%
% Macros for moduli stacks/spaces
%
\def\QCohstack{\mathcal{QC}\!\mathit{oh}}
\def\Cohstack{\mathcal{C}\!\mathit{oh}}
\def\Spacesstack{\mathcal{S}\!\mathit{paces}}
\def\Quotfunctor{\mathrm{Quot}}
\def\Hilbfunctor{\mathrm{Hilb}}
\def\Curvesstack{\mathcal{C}\!\mathit{urves}}
\def\Polarizedstack{\mathcal{P}\!\mathit{olarized}}
\def\Complexesstack{\mathcal{C}\!\mathit{omplexes}}
% \Pic is the operator that assigns to X its picard group, usage \Pic(X)
% \Picardstack_{X/B} denotes the Picard stack of X over B
% \Picardfunctor_{X/B} denotes the Picard functor of X over B
\def\Pic{\mathop{\mathrm{Pic}}\nolimits}
\def\Picardstack{\mathcal{P}\!\mathit{ic}}
\def\Picardfunctor{\mathrm{Pic}}
\def\Deformationcategory{\mathcal{D}\!\mathit{ef}}

\begin{document}
\title{Religion}
\maketitle
\tableofcontents

\section{Salvador field notes: Catholicism and Afro-Brazilian religions}
These notes began from visits in Salvador on September 26, 2026, especially
the Cathedral Basilica, MUNCAB, and the Casa do Benin. They are organized
around concrete questions and objects that called my attention.

\section{Catholic vocabulary encountered in Salvador}
\subsection{Jesuits}
The Jesuits are members of the Society of Jesus, a Catholic religious order
founded by Ignatius of Loyola and his companions and approved in 1540.
Education became one of its central activities. This connects the Jesuit
college in colonial Salvador with modern Jesuit institutions such as PUC--Rio
and the Universidad Iberoamericana in Mexico. Jesuits arrived in Bahia with
Tome de Sousa in 1549. The present Cathedral Basilica was built as the church
of the Jesuit college and only later became the cathedral. Cathedral names an
ecclesiastical function, the church containing the bishop's cathedra, rather
than an architectural size.

\subsection{One Mary, many invocations}
Nossa Senhora do Rosario, do Carmo, da Conceicao, das Dores, etc. are
different invocations or titles of Mary. Nossa Senhora do Carmo belongs to
the Carmelite tradition, named from Mount Carmel; the scapular is an important
Carmelite devotional object. The Immaculate Conception concerns Mary's own
conception: Catholic doctrine holds that Mary was preserved from original sin
from its first instant. It must not be confused with the virginal conception
of Jesus.

\section{A first map of Candomble}
Candomble is an Afro-Brazilian religion formed in Brazil from African
religious traditions carried through the Atlantic slave trade and reorganized
under slavery and diaspora. The orixas, voduns, and Central African religious
concepts did not originate in Brazil; the institutional formations called
Candomble did. African religion, Yoruba culture, and Candomble are therefore
not synonyms.

Nation in Candomble is not the same as a modern nation-state. A useful first
map is
\[
\begin{array}{c|c|c}
\text{tradition} & \text{major African matrix} & \text{divine beings}\\
\hline
\text{Ketu/Nago} & \text{Yoruba} & \text{orixas}\\
\text{Jeje} & \text{Ewe--Fon} & \text{voduns}\\
\text{Angola/Congo} & \text{Bantu/Central African} & \text{inquices/nkisi}.
\end{array}
\]
This is only a first map. The modern country Angola and the Candomble nation
Angola are different but historically related objects.

Nago is a Brazilian diasporic designation strongly associated with the Yoruba
world; Ketu is one important Candomble tradition within this broad field.
Ketu does not have one exclusive main orixa. Oxala is a major orixa associated
with creation, but a particular house can have its own patron. Olodumare or
Olorum is the supreme divinity and must not be confused with an orixa.
Likewise, \emph{Orun} is the spiritual realm, not the name Olorum.

\section{The Ile Axe Opo Afonja}
The Ile Axe Opo Afonja in Salvador is a historically important Ketu/Nago
terreiro, founded in 1910 by Eugenia Ana dos Santos (Mae Aninha), in a lineage
connected with the Casa Branca do Engenho Velho. Its patron is Xango Afonja.

An \emph{ialorixa} is the female sacerdotal leader of a terreiro; the
masculine counterpart is \emph{babalorixa}. An ialorixa leads a house, not
the whole Nago tradition.

The book \emph{Pedagogia da ancestralidade: fe, educacao e lideranca no Ile
Axe Opo Afonja} should be read as a deep study of one influential Ketu/Nago
house, not as a universal manual of Candomble. Its pedagogy of the terreiro
emphasizes coexistence, participation, hierarchy, oral memory, practice, and
the example of experienced members.

\subsection{Mae Stella de Oxossi}
Mae Stella de Oxossi (Maria Stella de Azevedo Santos, 1925--2018) was the
fifth ialorixa of the Opo Afonja, from 1976 until her death. The museum
timeline made her importance concrete: she worked for decades in nursing and
public health; visited Nigeria in 1981, including Oshogbo, Ile-Ife, and Ede;
participated in international conferences on orixa traditions; travelled to
Benin; became an author and public intellectual voice of Candomble; and in
2013 became the first Black woman and first ialorixa to enter the Academia de
Letras da Bahia. In the religious language encountered at the museum, her
death in 2018 was a return to \emph{Orun}.

\section{The 1983 manifesto and syncretism}
A museum panel reproduced the 1983 manifesto signed by Mae Stella de Oxossi,
Menininha do Gantois, Tete de Iansa, Olga de Alaketu, and Nicinha do Bogum.
This is an important primary source. The manifesto rejects the idea that
Candomble needs Catholicism for legitimacy. It treats syncretism as
historically connected with slavery and coercion, argues for practising
Candomble in its own terms, and criticizes folklorization, commercialization,
and touristic consumption. Historical saint--orixa associations are real, but
Candomble cannot be reduced to Catholic saints in disguise.

\section{Objects and concepts that called my attention}
\subsection{Ofa and Oxossi}
The bow-and-arrow form in Diogum's sculpture was recognizable as an
\emph{ofa}, an emblem of Oxossi. Diogum's name and work in iron point toward
Ogum, but this does not make every object he represents an emblem of Ogum.

\subsection{Beads}
Micangas are beads as material objects. In Afro-Brazilian religious contexts,
strings of beads can mark religious affiliations and may be ritually prepared.
Meaning depends on arrangement, colour, house, and ritual context. Three
museum encounters made the material visible: Hariel Revignet's bead works and
zigidas; Nadia Taquary's strings of beads in the colours of Iroko; and Sandro
Aiye's carved work using beads.

\subsection{Iroko and Iyami}
Nadia Taquary's \emph{Iroko: A arvore cosmica} connected Iroko with time,
ancestry, the gameleira, and the Iyami. The museum translation of Iyami as
feiticeiras should be treated cautiously: the Yoruba conception of ancestral
feminine power is richer than the European category witch.

\subsection{Saint George}
Saint George is a Christian warrior saint. Afro-Brazilian saint--orixa
associations vary regionally and between houses; they are not universal
equivalences. The plant \emph{lanca-de-Sao-Jorge} adds another layer: a
popular protective plant whose name evokes the saint's weapon.

\section{Afro-Atlantic comparison}
Candomble Ketu/Nago and Cuban Santeria are distinct American religious
formations that preserve and transform Yoruba inheritances. Santeria also
cultivates orishas. This is useful for thinking about the Yoruba diaspora
without collapsing the religions into one.

The Casa do Benin added the geographical side: Benin, Nigeria, Bahia, and the
Atlantic are concrete places linked by forced migration, religious
transmission, return journeys, scholarship, and cultural exchange. Pierre
Verger's idea of flux and reflux gives a useful question: what travelled from
Africa to Bahia, what changed in Bahia, and what later travelled back?

\section{Reading leads}
\emph{Pedagogia da ancestralidade} is a study of Opo Afonja and its
pedagogy/leadership. \emph{Nacoes do Candomble: Jeje}, organized by Mateus
Aleluia Lima, is valuable precisely because it moves outside the Ketu/Nago/
Yoruba emphasis of much of the day's material and studies a Jeje matrix. It
should not by itself be treated as a general manual of all Candomble.

\bibliography{my}
\bibliographystyle{amsalpha}
\end{document}
